\section{Discussion and Limitations}

\subsection{Why Protocol Split Matters}

The single-goal mission protocol and the goal-reissue stress protocol expose
different method rankings. Wind-Level 0.85 is strongest in the
single-goal setting, while Fixed Scale 0.80 is strongest under goal-reissue
stress. A single combined aggregate would hide this structure and could make a
method appear broadly robust when it is specialized to one regime.

\subsection{Why Strong Simple Baselines Matter}

The strong heuristic and static baselines are not a weakness of the evaluation.
They reveal that execution governance is a real operating-point problem:
simple scales can work well when matched to a protocol, but they also expose
trade-offs across protocols. Their strength makes the balanced robustness
claim for Risk Adapter v1 more meaningful because it is measured against strong
specialists rather than weak baselines.

\subsection{Balanced Robustness as a Protocol-Level Criterion}

Risk Adapter v1 is the balanced protagonist because it has the highest mean
valid count, highest worst-protocol count, and lowest total regret under the
tested protocols. This does not imply statistical superiority, universal
dominance, or guaranteed safety. The narrower claim is that, among the
evaluated learned / risk-conditioned governors, it gives the most balanced
protocol-level behavior.

\subsection{Failure-Mode-Aware Interpretation}

Invalid-only failure groups help interpret where failures concentrate, using
categories such as command/control upstream, planner/reference upstream, and
state/task upstream. These labels are diagnostic rather than definitive
root-cause attribution. Representative traces support mechanism hypotheses,
including early pre-arrival vulnerability in Trial 4 stress failures and
delayed post-arrival divergence in a Trial 6 single-goal bottleneck.

\subsection{Why a New Variant Is Not Introduced Yet}

The representative trace review did not justify a simple threshold patch.
Trial 4 stress weaknesses include early pre-arrival vulnerability and risk
availability or timing issues, while Trial 6 shows that lower scale is not
automatically safer. A future governor should be phase-aware, reference-aware,
risk-health-aware, or failure-aware only if a reusable mechanism is established.
Trial-specific patching is outside the current paper claim.

\subsection{Limitations}

The evidence is simulation-only and should not be described as hardware or
field robustness evidence. The repeated-run counts are descriptive and do not
establish statistical significance. The execution governor is empirical and
does not provide a formal safety guarantee. Failure-group labels are diagnostic
and do not establish definitive physical causality. The main evaluation focuses
on Trials 4, 5, and 6, and Risk Adapter v2 lacks a matched stress result for
the balanced table. The source, training, calibration, and confidence behavior
of the risk scores also require clearer documentation before a stronger journal
extension.

\subsection{Future Work}

Future work should investigate trace-driven phase-aware governors, risk
availability or confidence fallback, broader fixed-scale and wind-level
frontier sweeps, risk calibration and delayed/shuffled risk ablations,
speed-only versus acceleration-only scaling, additional wind profiles, payload
masses, cable lengths, and hardware or HIL validation if available.
